\documentclass[12pt, letterpaper]{article}

%-------Packages---------
\usepackage{amssymb,amsfonts}
\usepackage{mathptmx}
\usepackage{amsthm}
\usepackage{graphicx}
\usepackage[all,arc]{xy}
\usepackage[colorlinks=true, allcolors=blue]{hyperref}
\usepackage{enumerate}
\usepackage{bbm}
\usepackage{dsfont}
\usepackage{mathrsfs}
\usepackage{tikz-cd}
\usepackage{setspace}
\usepackage{import}
\usepackage[utf8]{inputenc}
\usepackage{hyperref}
\usepackage{tikz}
\usepackage{float}
\usepackage{mdframed}
\usepackage[nocheck]{fancyhdr}
\usepackage[nointegrals]{wasysym}

\usepackage[left=1in,top=1in,right=1in,bottom=1in]{geometry}

\usepackage{mathtools}
\usepackage{setspace}
\usepackage{cleveref}
\usepackage{multicol}
\usepackage{mathpazo}

\usepackage{algorithm}
\usepackage[noend]{algpseudocode}
\usetikzlibrary{positioning}

\newcounter{problem}

% Define a new command for problems
\newcommand{\q}{
    \newpage % Start a new page
    \stepcounter{problem} % Increment problem counter
    \subsection*{Problem \arabic{problem}} % Display the problem counter
}
\newcommand{\C}{\mathbb{C}}
\newcommand{\R}{\mathbb{R}}
\newcommand{\Q}{\mathbb{Q}}
\newcommand{\Z}{\mathbb{Z}}
\newcommand{\N}{\mathbb{N}}
\renewcommand{\P}{\mathbb{P}}
\newcommand{\F}{\mathbb{F}}
\newcommand{\contr}{\Rightarrow \Leftarrow}
\newcommand{\HRule}[1]{\rule{\linewidth}{#1}}
\newcommand{\s}{\(\,\)}
\renewcommand{\t}[1]{\text{#1}}
\newcommand{\ang}[1]{\left\langle #1 \right\rangle}

% Meta data
\setstretch{1.2}

\begin{document}
\setlength{\parindent}{0pt}
\pagestyle{fancy}
\fancyhf{}
\fancyhead[LOE]{\slshape{\textbf{Vincent Ferrara}}}
\fancyhead[ROE]{\thepage}

\q
\begin{enumerate}[(a)]
    \item Base Case:
    
    For $n=1$ it is trivial since $I \subseteq p_1$.

    Inductive Step: Assume true for $2 \leq n-1$. Consider $n$.
    
    If $I$ is a subset of any $n-1$ unions of primes then we are done thus we can assume
    \[I \nsubseteq \bigcup_{j \neq i}p_j\]
    for all $i \in \{1,\dots,n\}$
    
    Thus let
    \[x_i \in I \setminus \bigcup_{j \neq i}p_j\]

    If $x_i \in p_i$ since $I \subseteq p_1 \dots p_2$

    Consider $x=x_1+x_2x_3\dots x_n$

    If $x \in p_i$ for $1 < i \leq n$ then $x-x_2x_3\dots x_n = x_1 \in p_i$ since $x_2x_3\dots x_n \in p_2 \cap p_3 \cap \dots \cap p_n$.

    This is a contraction since $x_1 \notin p_2 \cup p_3 \dots \cup p_n$.

    If $x \in p_1$ then $x-x_1=x_2x_3\dots x_n \in p_1$ this is a contraction since $p_1$ is a prime ideal this means that some $x_i \in p_1$ for $i \neq 1$.

    This $x \notin p_1 \cup p_2 \cup \dots \cup p_n$ which is a contraction thus
    \[I \subset \bigcup_{j \neq i}p_j\]
    for some $i \in \{1,\dots,n\}$
    Thus by the Inductive Hypothesis we are done.
\end{enumerate}

\q
\begin{enumerate}[(a)]
    \item Let $\varphi: A \to A[x] \diagup (fX-1)$ s.t. $\varphi(a)=[a]$.
    
    Consider $\varphi(f^k)$ for $k \in \N\cup \{0\}$.

    $\varphi(f^k)=\varphi(f)^k=[f]^k$.

    Consider $[f]^kX^k=([f]X)^k=1$ since $fX-1=0 \implies fX=1$.

    Thus, $f^k$ is a unit.

    Let $g(x) \in A[X] \diagup (fX-1)$. Thus $g(x)=a_0+a_1X+\dots+a_nX^n$ for $a_i \in A$

    since $X=\varphi(f)^{-1}$ we get that 
    \[g(X)=a_0+a_1+f^{-1}+\dots+a_nf^{-n}\]
    Now combining terms we get
    \[g(X)=\dfrac{a_0f^n+a_1f^{n-1}+\dots+a_n}{f^{n}}\]
    Thus
    \[g(X)=\varphi(a')\varphi(f^n)^{-1}\]
    Thus every thing is written in the form $\varphi(a)\varphi(s)^{-1}$ for $s \in \{1,f,f^2,\dots\}$

    Let $b \in \ker{\varphi}$ thus $\varphi(b)=(fX-1)g(X)$ for some $g(X)=g_0+g_1X+\dots+g_nX^n \in A[X]$
    \[b=-g_0+(fg_0-g_1)X+(fg_1-g_2)X^2+\dots+(fg_{n-1}-g_n)X^n+fg_nX^{n+1}\]

    \begin{itemize}
        \item (Constant)\[g_0=-b\]
        \item ($X$) $fg_0-g_1=0$ \[g_1=fg_0=-fb\]
        \item ($X^i$) $fg_{i-1}-g_{i}=0$ \[g_i=fg_{i-1}=fg_{i-1}=f^{i}g_0=-f^ib\]
        \item ($X^{n+1}$) $fg_n=0$ \[fg_n=-f^{n+1}b=0\]
        Thus
        \[f^{n+1}b=0\]
    \end{itemize}
    Thus $\ker{\varphi} \subseteq \{a \in A \mid \exists k \geq 0,\; f^k[a]=0\}$

    Let $b \in \{a \in A \mid \exists k \geq 0,\; f^ka=0\}$ thus $f^kb=0$ for $k \geq 0$

    Consider $-b(1+fX+f^2X^2+\dots+f)(fX-1)$
    \begin{align*}
        -b(1+fX+f^2X^2+\dots+f^{k-1}X^{k-1})(fX-1)=-b((fX)^k-1) \intertext{finite geometric series}
        =b-bf^kX^k=b
    \end{align*}
    Thus $b \in (fX-1)$ thus $\varphi(b)=0$ thus $b \in \ker{\varphi}$.

    Thus, there is a unique isomorphism between $A_f$ and $A[X] \diagup (fX-1)$

    Therefore,
    \[\Z[X] \diagup (2X-1) \cong \{1,2,4,\dots\}^{-1}\Z=A\]
    
    Let $\varphi: A \to \Q$ s.t. $\varphi(a/2k)=a/2k$

    Injective: This is clearly injective since $a/2k = b/2m$ in $\Q$ then $a2m=b2k \implies 1(a2m-b2k)=0$ which satisfies the equivalence relation in $A$.

    Surjective:

    Let $a/b \in \Q$ then $a/b=2a/2b$. Since $2b \in \{1,2,4,\dots\}$, we get that $2a/2b \in A$. Thus $\varphi(2a/2b)=a/b$.

    Homomorphism:

    $\varphi(a/b)\varphi(c/d)=(a/b)(c/d)=ac/bd=ac/bd=\varphi(ac/bc)$

    $\varphi(a/b)+\varphi(c/d)=(ad+cb)/bd=\varphi((ad+cb)/bd)=\varphi(a/b+c/d)$

    $\varphi(1/1)=1/1$

    Thus $A \cong \Q$

    \[(\Z \diagup 6 \Z)[X] \diagup(2X-1) \cong \{1,2,4\}^{-1}(\Z \diagup 6\Z) \cong \Z \diagup 3\Z\]
    We showed the last part in class

    \item Let $\varphi$ be the natural homomorphism from $A \to S^{-1}A$
    
    Injective: Let $a \in \ker(\varphi)$. Thus, $\varphi(a)=0$. Thus, by definition of $S^{-1}A$ there is a $u \in S$ s.t. $ua=0$ since $u \in S$ it is a unit thus we can multiply its inverse on both sides and get $a=0$.

    Thus injective

    Surjective: The map is by definition surjective.

    Homomorphism:

    Proven in class it is a homomorphism

    Thus $A \cong S^{-1}A$

    If $K$ is a field then $S^{-1}K$ where $S$ is all units is actually just $(K\setminus \{0\})^{-1}K$ since all things in $K$ are units except zero. Thus, by definition this is the same as $Frac(K)$ thus by the proof above $Frac(K) \cong K$

    Consider $Frac(R)=(R \setminus \{0\})^{-1}R$

    \begin{itemize}
        \item We get Associativity, Commutativity, Additive/Multiplicative Inverse, Additive inverses, Distributivity for free.
        \item (Multiplicative Inverses) Let $x \in Frac(R)$ s.t. $x \neq 0$. Thus, $x=\varphi(a)\varphi(b)^{-1}$ for $a,b \in R \setminus \{0\}$
        
        Consider $y = \varphi(b)\varphi(a)^{-1}$ this is well-defined since $a \neq 0$.

        Thus $xy=\varphi(a)\varphi(a)^{-1}\varphi(b)\varphi(b)^{-1}=\varphi(1)\varphi(1)^{-1}=1$.

        \item ($0 \neq 1$) Since two non zero things in $R$ never multiply to zero we get that there is no element $s \in R \setminus \{0\}$ s.t. $s(1-0)=0$. 
    \end{itemize}
    Thus, $Frac(R)$ is a field thus $Frac(Frac(R))=Frac(R)$

    \addtocounter{enumi}{2}\item Seperate Proof: Show $S^{-1}R \cong S'^{-1}R$ where $S'=(S \setminus{\t{units of } S}) \cup \{1\}$ for non zero rings.
    
    Assume that $S'$ has more than just 1 since if it has just one then $S$ is all units, so then you can look at (b).

   \begin{itemize}
    \item  Show $S'$ is multiplicative set with $1$
 
     It has $1$ by definition.
 
     Let $a,b \in S' \setminus {1}$ then since $a,b \in S \implies ab \in S$

     Assume $ab$ is a unit in $R$ with inverse $c \in R$ this means that
     \[(ab)c=1 \implies a(bc)=1\]
     Thus $a$ is a unit which is a contraction. Thus, $ab \in S'$

     Also multiplicative with 1 is trivial.
   \end{itemize}

   Let $r/s \in S^{-1}R$

   If $s \in S'$ then there is nothing to do.

   If $s$ is a unit, then choose any element $t \in S' \setminus \{1\}$ thus $st \in S$.

   Assume $st$ is also a unit then $(st)c=1 \implies s(tc)=1 \implies t=s^{-1}c^{-1}$ this $t$ is a unit which is a contraction.

   Thus $st \in S'$ now we can represent $r/s=rt/st$

   Thus every element in $S^{-1}R$ can be written in the form of $r/s'$ for $s' \in S'$

   Now using the classifications as above take the natural map $\varphi:(S'^{-1})R \to S^{-1}R$ s.t. $\varphi(r/s')=r/s'$.

   This is well defined for the reasons above.
   
   This is a homomorphism since we're using the same multiplication and addition rules.

   This is surjective for the reasons above.

   Injective:

   Assume consider $r/s',r'/s'' \in S^{-1}R$ s.t. $\varphi(r/s')=\varphi(r'/s'')$ thus $r/s'=r'/s''$ in $S^{-1}R$ thus there is an element in $t \in S$ s.t.
   \[t(rs''-r's')=0\]
   If $t \in S'$ then we are done.

   If $t$ is not in that set then we know it is a unit thus we can get $rs''-r's'=0$ since $1$ is in $S'$ we know thus means that these elements are equivalent in $S'^{-1}R$

   Thus injective.

   Thus $S^{-1}R \cong S'^{-1}R$

   Now for $\Z \diagup 60\Z$ using the proof above all the units are things that are co prime with $60$. Thus, we only have to look at things that divide 60, which are in the form $2^a3^b5^c$ where $a,b,c \geq 0$.

   However, for a localization with $2^a3^b5^c$ for $a,b,c \geq 1$. This gives us an inverse of $2,3,5$. Since $1/(2^a3^b5^c) \times 2^a3^b5^{c-1}/1 \equiv 1/5$. Thus this problem actually reduces to just choosing subsets of $\{2,3,5\}$ in your set $S$.

   Thus

   \begin{enumerate}[1.]
    \item $(S=\{1\})$ This just gives us $R$ by part above.
    
    \item $(2 \in S)$ Just $2$ in our set. Which is $\Z \diagup 15 \Z$
    
    \item $(3 \in S)$ Just $3$ in our set. Which is $\Z \diagup 20 \Z$
    
    \item $(5 \in S)$ Just $5$ in our set. Which is $\Z \diagup 12 \Z$
    
    \item $(2,3 \in S)$ Just $2,3$ in our set. Which is $\Z \diagup 5 \Z$
    
    \item $(2,5 \in S)$ Just $2,5$ in our set. Which is $\Z \diagup 3 \Z$
    
    \item $(3,5 \in S)$ Just $3,5$ in our set. Which is $\Z \diagup 4 \Z$
    
    \item $(2,3,5 \in S)$ This gives us that everything in our set has an inverse, so for example $2*30=0 \implies 1=0$ since $2,30$ have an inverse. Thus, we have the zero ring in this case.
   \end{enumerate}

    \item Let $P$ be a prime ideal in $\Z[X]$. Since $P \cap \Z$ is a prime ideal in $\Z$ we get that $P \cap \Z =(0),(p)$ for $p$ prime.
    
    If $\Z \cap P = (0)$, then no integer is in $P$. Thus consider $Frac(\Z)[X]$ since $P$ does not meet $\Z \cap \{0\}$ this means its extension is a prime ideal in $Frac(\Z)[X]=\Q[X]$

    Since $\Q$ is a field we know that $\Q[X]$ is a $PID$ thus $P^e=(f(X))$ for some irreducible polynomial in $\Q[X]$. Now let $g(X) \in \Z[X]$ s.t. $g(X)/c = f(X)$ for $c \in \Z$. Thus $(g(X))=(f(X))$, also we can cancel the coefficients, so we get that their $\gcd$ is 1.

    Separate: 

    Let $f(X) \in P$ then $f(X)/1 \in P^e \cap \Z[X]$

    Let $f(X) \in P^e \cap \Z[X] \implies f(X)=h(X)/s$ for $h(X) \in \Z[X]$ and $s \in \Z \setminus \{0\}$ $\implies sf(X)-h(X)=0 \implies sf(X)=h(X) \in P$ since $s \notin P \implies f(X) \in P$ thus $P^e \cap \Z[X]=P$.

    Consider $(g(X)) \subseteq \Z[X]$. Let $h(X) \in (g(X))$ thus $h(X)=g(X)q(X)$ for $q(X) \in \Z[X]$.

    Now we can extend this to $\Q[X]$ just by dividing by 1, and we get $h(X)=g(X)q(X)=cf(X)q(X) \in \Q[X] \implies h(X) \in P^e \cap \Z[X]=P$

    Let $h(X) \in P$ then extend this to $\Q[X]$ thus $h(X) \in P^e=(f(X))=(g(X))$ thus $h(X)=g(X)q(X)$ for $q(X) \in \Q[X]$. Now get rid of all the denominators in $q(X)$ and we get that $zh(X)=g(X)q_1(X)$ for $q_1(X) \in \Z[X]$. Thus $zh(X) \in P^e \cap \Z[X]=P$, since $P \cap \Z =0 \implies h(X) \in P$.

    Thus, $P=(g(X))$ which is an irreducible polynomial in $\Z[X]$.

    If $\Z \cap P = (p)$. Then consider the projection from $\Z[X] \to \F_p[X]$. Since there is a bijection between ideals in $\Z[X]$ that contain $(p)$ and ideals in $\F_p[X]$. Thus $P \diagup (p)$ is a prime ideal in $\F_p[X]$.

    Since $\F_p$ is a field we know that the ideals are principle in $\Z[X]$ thus the prime ideals are $(0)$ and $(f(X))$ for irreducible in $\F_p$

    If $P\diagup(p)=(0)$ then $P=(p)$.

    If $P \diagup (p) \neq (p)$ then it will equal $P \diagup (p) = (f(X))$ where $f(X)$ is an irreducible polynomial mod $p$. Thus, $P=(p,f(X))$


    \addtocounter{enumi}{2} \item 
    Ring
    
    For two germs $[f]$ and $[g]$ in $\mathbb{R}_x$, define:
    \[
    [f] + [g] := [f+g] \quad \text{and} \quad [f] \cdot [g] := [f \cdot g].
    \]
    Here, $f+g$ and $f\cdot g$ are defined on the intersection of the domains of definition of $f$ and $g$. 
    
    Well-Defined: Suppose $[f]=[f']$ and $[g]=[g']$, meaning there exist open neighborhoods $W_1$ and $W_2$ of $x$ such that
    \[
    f|_{W_1} = f'|_{W_1} \quad \text{and} \quad g|_{W_2} = g'|_{W_2}.
    \]
    Then on $W = W_1\cap W_2$, we have
    \[
    (f+g)|_W = f|_W + g|_W = f'|_W + g'|_W = (f'+g')|_W,
    \]
    and similarly for the product. Thus, the operations are independent of the choice of representatives, making $\mathbb{R}_x$ a ring with additive identity $[0]$ and multiplicative identity $[1]$.
    
    
    Local Ring
    
    Let
    \[
    \mathfrak{m}_x = \{ [f] \in \mathbb{R}_x \mid f(x)=0 \}.
    \]
    If $[f], [g] \in \mathfrak{m}_x$, then
    \[
    (f+g)(x)=f(x)+g(x)=0,
    \]
    so $[f]+[g] \in \mathfrak{m}_x$. Also, for any $[h]\in \mathbb{R}_x$,
    \[
    (f\cdot h)(x)=f(x)h(x)=0,
    \]
    so $[f]\cdot[h] \in \mathfrak{m}_x$.
    
    Let $[f]\in \mathbb{R}_x$ with $[f]\notin \mathfrak{m}_x$. Then $f(x)\neq 0$. By the continuity of $f$, there exists an open neighborhood $U$ of $x$ where $f(y)\neq 0$ for all $y\in U$. Define
    \[
    g(y)=\frac{1}{f(y)} \quad \text{for } y\in U.
    \]
    Since $g$ is smooth on $U$, we have
    \[
    [f]\cdot [g] = [f\cdot g] = [1].
    \]
    Thus, $[f]$ is invertible.
    
    \vspace{0.5em}
    
    Since every element not in $\mathfrak{m}_x$ is a unit, the complement of $\mathfrak{m}_x$ consists entirely of units, and hence $\mathfrak{m}_x$ is the unique maximal ideal of $\mathbb{R}_x$. Therefore, $\mathbb{R}_x$ is a local ring.    

    \item Let $x \in \mathfrak{N}(A)$ thus $x^n=0$ for $n \in \N$ but since $0 \in p$ for all $p \in Spec(A)$ this implies by the definition of a prime ideal that $x \in p$ since for $x^n=0 \in p$ some factor of $x^n$ has to be in $p$ but all the factors are $x$ thus $x \in p$ for all $p \in Spec(A)$
    
    Let $x \in p$ for all $p \in Spec(A)$. Assume for contraction that $x$ is not nilpotent thus consider $A_x$. Since $0$ was not in our multiplicative set we know that $A_x$ is not the zero ring thus there exists a maximal ideal $m$ thus prime. However this is a contradiction since there is a bijection from primes ideals in $A_x$ to ideals prime ideals of $A$ that do not meet $\{1,x,x^2,\dots\}$. Since $x \in p$ for all $p \in Spec(A)$. This means that $x$ must be nilpotent.

    Therefore, the intersection of all prime ideals is the nilradical.

    Let $I$ be an ideal in $A$. Let $P$ be the intersection of all prime ideals of $A$ containing $I$.

    Let $x \in r(I)$ thus $x^n \in I$ for $n \in \N$. Thus, $x \in P$ by the same reason as above.

    Let $x \notin r(I)$ thus no power of $x$ is in $I$. Consider $A \diagup I$ s.t. the image of $x$ is $y$ consider the set $S=\{1,y,y^2,\dots\}$. Since $x \notin r(I)$ this means that $y$ will not be nilpotent. Thus $(A \setminus I)_{y}$ is non zero thus we can get a maximal ideal $m$ which is prime. Thus, this corresponds to a prime ideal in $A \diagup I$ that does not meet $S$. Now contract this prime in $A \diagup I$ to $A$ and we get a prime ideal in $A$ that contains $I$ but not $x$. Since if it did have $x$ then we can take its image then the prime ideal in $A \diagup I$ would have $y$ which is a contraction of it not having $S$. Thus, $x \notin P$.

    Therefore $r(I)=P$
    
    \item If $A=(0)$ then $Spec(0)=\emptyset$, so it is Hausdorff.
    
    Assume $Spec(A)$ where $A \neq (0)$ is Hausdorff. Assume for contradiction there is a $p \in Spec(A)$ that is not maximal. Thus, there is a maximal ideal $q$ s.t. $p \subsetneq q$. Thus, there are open sets $p \in U$ and $q \in V$ s.t. $U \cap V= \emptyset$. Using the basis for this space we get that there is an $f \in A$ s.t. $q \in X_f \subseteq V$.

    Since $p \subseteq q \implies f \notin p$. Thus, $p \in X_f$ which is a contraction that $U \cap V =\emptyset$. Thus all prime ideals are maximal.

    Assume every prime ideal of $A$ is maximal. Let $p,q \in Spec(A)$. Since $p,q$ are maximal you can find elements $f \in p$ and $f \notin q$.
    
    Consider $S=A \setminus q$ for $S^{-1}A$. Since all prime ideals in $S^{-1}A$ are prime ideals in $A$ that do not meet $S$ we get that $q^e$ is the unique prime ideal. Thus, for $g/1 \in q^e=r(0)$ by the last parts. Thus, $g^n/1=0/1$. Thus, there is an $a \in S$ s.t. $ag^n=0$. Since $0 \in p \cap q$ and $g^n \notin p$ thus $a \in p$. Thus $p \in X_a$ and $p \notin X_f=X_{f_n}$ same for $q$.

    By (a) we know that $X_{f^n} \cap X_{a} = X_{f^na}=X_{0}=\emptyset$. Thus, $Spec(A)$ is Hausdorff.
\end{enumerate}

\q
\begin{enumerate}[(a)]\addtocounter{enumi}{1}
    \item Assume $f$ is a unit. Consider $X_f=\{p \in X \mid f \notin p\}$. Let $fg=1$. If $f \in p \implies fg=1 \in p$, but since $p$ is proper ideal this means $f \notin p$. Thus $X_f=X$
    
    Assume $X=X_f$. Thus, $f \notin p$ for all $p \in X$. Thus, $f$ is not in any maximal ideal, since all maximal ideals are prime. Consider $(f)$. Assume for contraction it is a proper ideal. By Zorn's Lemma we know that $(f) \subseteq (m)$ where $(m)$ is maximal but $f$ is not in any maximal ideal thus $(f)=A$. Thus $fa=1$ for $a \in A \implies f$ is a unit.

    \item Assume $f$ is nilpotent. Thus $f^k=0$ for $k \in \N$. Since $0 \in p$ for all $p \in X$. This means that $f^k \in p$. By definition of a prime ideal the only way for $f^k \in p$ is if one of its factors is in $p$, but since all of its factors are $f$ this means $f \in p$.
    
    Assume $X_f=\emptyset$ Thus $f$ is in the intersection of all the prime ideals which shown in problem $2$ is the nilradical. Thus, $f$ is nilpotent.

    \addtocounter{enumi}{1}\item Proven in last homework that those basic open sets form a basis thus we only have to consider open covers of the form $\{X_{f_i}\}_{i \in I}$ since every open set can be written as a union of basis sets.

    So let $\{X_{f_i}\}_{i \in I}$ be an open cover for $X$.

    Thus,
    \[X=\bigcup_{i \in I}X_{f_i}=\bigcup_{i \in I}(X \setminus V(f_i))=X \setminus \bigcap_{i \in I}V(f_i)\]
    Shown in last homework this becomes
    \[X=X \setminus \bigcap_{i \in I}V(f_i)=X \setminus V(\bigcup_{i \in I} \{f_i\})=X \setminus V((f_i \mid i \in I))\]
    Since $X=X \setminus V((f_i \mid i \in I))$. This means that $V((f_i \mid i \in I))=\emptyset$. If the ideal $(f_i \mid i \in I)$ is proper, then by Zorn's lemma it is contained in a maximal ideal $m$, but maximal ideals are prime so $m \in V((f_i \mid i \in I))$ which is a contraction thus $(f_i \mid i \in I)=A$.

    This means there exists a finite combination of $f_i$ and $a_i \in A$ s.t. $1=a_1f_{i_1}+a_2f_{i_2}+\dots+a_nf_{i_n}$

    Let $p \in X$. Suppose $p \notin X_{f_{i_1}} \cup \dots \cup X_{f_{i_n}}$

    Thus, $f_{i_j} \in p$ for all $j \in \{1,\dots,n\}$ thus $a_1f_{i_1}+a_2f_{i_2}+\dots+a_nf_{i_n}=1 \in p$ which is a contraction thus $p \in X_{f_{i_1}} \cup \dots \cup X_{f_{i_n}}$

    Thus $X \subseteq X_{f_{i_1}} \cup \dots \cup X_{f_{i_n}}$

    Therefore, $X=X_{f_{i_1}} \cup \dots \cup X_{f_{i_n}}$ thus $X$ is quasi compact.
\end{enumerate}
\end{document}